\par\noindent\textbf{(1)} 正负离子局域密度 n\_±=n₀ exp(∓eφ/k\_BT)，电荷密度 e(n\_+−n\_−)≈−2n₀e²φ/(k\_BT)。Poisson式 −ε∇²φ=Qδ³(r)+ρ\_ion，重排 (∇²−κ²)φ=−Qδ³(r)/ε，其中 κ²=2n₀e²/(εk\_BT)。
\par\noindent\textbf{(2)} 三维Helmholtz算符衰减Green函数为 e\textasciicircum{}\{−κr\}/(4πr)，故 φ(r)=Qe\textasciicircum{}\{−κr\}/(4πεr)。由ρ\_ion=−εκ²φ计算∫ρ\_ion d³r=−Qκ²∫₀∞r e\textasciicircum{}\{−κr\}dr=−Q，说明完整离子云恰补偿中心电荷，远处总电通量趋零。
\par\noindent\textbf{(3)} 径向电场 E\_r=−φ'=Qe\textasciicircum{}\{−κr\}(1+κr)/(4πεr²)。点电荷本征静电能在r→0仍发散，不能宣称屏蔽后全部场能有限；以相同芯截止比较屏蔽与未屏蔽点电荷的差，线性响应自能修正为 ΔU=Qφ\_ind(0)/2=−Q²κ/(8πε)，其中 φ\_ind(0)=−Qκ/(4πε)。
\par\noindent\textbf{(4)} 线化前提是主要贡献区域 |eφ|≪k\_BT，故在试验电荷近旁通常失败，须设有限半径芯a并限制 Qe/(4πεa k\_BT)≪1。κ→0 还原 φ=Q/(4πεr)，诱导云总电荷的极限应按无穷体积与κ顺序理解，否则会误把逐点趋零的电荷密度与总积分混淆。
